\documentclass[12pt]{ctexart}

\usepackage[margin=1.8cm]{geometry}
\usepackage{amsmath}
\usepackage{booktabs}
\usepackage{tabularx}
\usepackage{enumitem}
\usepackage{titlesec}
\usepackage{xcolor}
\usepackage[hidelinks]{hyperref}

\emergencystretch=2em

% 中文字体：使用系统已安装的思源字体（西文走 Liberation Serif）
\usepackage{xeCJK}
\setmainfont{Liberation Serif}
\setCJKmainfont{Source Han Serif SC}
\setCJKsansfont{Source Han Sans SC}
\setCJKmonofont{Source Han Mono SC}

% --- minimal styling ---
\definecolor{primary}{RGB}{15,15,35}
\titleformat{\section}{\Large\bfseries\color{primary}}{}{0em}{}
\titlespacing*{\section}{0pt}{0.6em}{0.2em}
\titleformat{\subsection}{\large\bfseries\color{primary}}{}{0em}{}
\titlespacing*{\subsection}{0pt}{0.6em}{0.15em}
\setlist{nosep, topsep=0.25em}

\setlength{\parskip}{0.15em}
\setlength{\parindent}{0pt}
\linespread{1.1}

\newcommand{\coursecode}{2026 Fall}
\newcommand{\coursename}{AI Agent: Practice}

\begin{document}

\pagestyle{empty}

\begin{center}
  {\LARGE\bfseries \coursecode{} -- \coursename}\\[0.4em]
  {\small Instructor: 林逸涵 \texttt{yhlin@eitech.edu.cn}}
\end{center}
\vspace{0.25em}

\section{Overview}

Agents are software systems centered on Large Language Models; they interact with their environment---invoking tools, writing code, and retrieving information---through perception, reasoning, planning, and action to accomplish multi-step, real-world tasks.

The ultimate goal of this seminar is to enable students to understand the working principles of agents and confidently select the agents that meet their needs.

\section{Goal}

By the end of the course, students should be able to:

\begin{itemize}
  \item Explain the fundamental components of an intelligent agent (the ReAct loop) and mainstream design paradigms;
  \item Configure tool calling (or tool use) for an LLM and design appropriate tool interfaces;
  \item Build agents capable of retaining context by leveraging context management, Retrieval-Augmented Generation (RAG), and memory mechanisms;
  \item Decompose complex tasks into multi-agent collaborative workflows and implement them using mainstream frameworks;
  \item Design evaluation metrics for their agents and conduct demonstrations.
\end{itemize}

\section{Schedule}

\begin{center}
\begin{tabularx}{0.8\textwidth}{@{}c c >{\raggedright\arraybackslash}X@{}}
  \toprule
  \bfseries Week & \bfseries Date & \bfseries Topic \\
  \midrule
  1 & Sep 11  & LLM; Agent \& the ReAct Loop; Tool Calling Basics; Harness \\
  2 & Sep 18  & Tool Design in Depth \\
  3 & Sep 25  & Context, RAG \& Skills \\
  4 & Oct 2   & Planning \& Multi-Agent \\
  5 & Oct 9   & Evaluation \& Wrap-up \\
  \bottomrule
\end{tabularx}
\end{center}

\noindent\textbf{Time:} Every Friday 14:00--16:00 \quad \textbf{Location:} 郑坚利红楼 413

\noindent 5 sessions, 1h lecture + 1h hands-on each

\section{Reading}

\begin{description}[leftmargin=6em, font=\normalfont\emph]
  \item[ReAct] Synergizing Reasoning and Acting in Language Models \href{https://arxiv.org/abs/2210.03629}{(arXiv:2210.03629)}
  \item[CodeAct] Code as Actions for Language Model Agents \href{https://arxiv.org/abs/2402.01030}{(arXiv:2402.01030)}
  \item[mini-swe-agent] A minimal coding agent \href{https://github.com/swe-agent/mini-swe-agent}{(GitHub: swe-agent/mini-swe-agent)}
\end{description}

\section{AI Policy}

The course \emph{encourages} the use of agent-based programming tools to assist with the hands-on sessions, and requires that any tools used be noted in the submission.

\end{document}
